A coleta ocorreu em 21 de outubro de 2026, no Windows 10, Python 3.14.7 e AMD Radeon Vega 8, com OpenGL 3.3. Foram medidos 4.320 quadros (540 por combinação) após descartar 1.440 de aquecimento. O destino offscreen de C4 manteve $1920\times1080$ pixels, apesar dos limites do monitor.

\begin{table}[t]

\centering

\caption{Resultados agregados de três repetições.}

\label{tab:resultados}

\begin{tabular}{llrrr}

\toprule

Config. & Modo & Média (ms) & P99 (ms) & FPS \\

\midrule

C1 & Mono & 20.83 & 28.38 & 48.01 \\

C1 & Multi & 23.19 & 31.71 & 43.13 \\

C2 & Mono & 21.30 & 30.68 & 46.94 \\

C2 & Multi & 23.81 & 35.09 & 42.00 \\

C3 & Mono & 20.87 & 27.57 & 47.91 \\

C3 & Multi & 24.16 & 27.74 & 41.40 \\

C4 & Mono & 30.76 & 47.63 & 32.51 \\

C4 & Multi & 34.43 & 45.08 & 29.05 \\

\bottomrule

\end{tabular}

\end{table}

A média multicamada foi maior nas quatro configurações: 11.3\%, 11.8\%, 15.7\%, 11.9\%, respectivamente. Esse resultado é compatível com o custo adicional de destinos intermediários e passes de tela cheia. Ele não isola o tempo de cada mecanismo, pois a medida conjunta inclui CPU, rasterização e sincronização.

Em C4, as médias foram 30,76 ms (monolítica) e 34,43 ms (multicamada), correspondendo a 32,51 e 29,05 FPS. O P99 multicamada, 45,08 ms, foi inferior aos 47,63 ms monolíticos nessa coleta. Portanto, média e cauda não devem ser confundidas, nem se pode concluir superioridade temporal geral por uma única métrica.

Todos os quadros medidos ultrapassaram o orçamento de 16,67 ms. A cena desta avaliação não sustentou 60 FPS sincronizados nessa máquina. A fluidez interativa observada em outras condições não altera esse resultado do experimento.

Os valores de C1 a C3 permaneceram próximos e não cresceram monotonicamente com a carga. O custo da cena 3D, a posição concentrada dos retângulos e a variação do sistema limitam inferências sobre escalabilidade. C4 altera simultaneamente resolução e carga; esse contraste não permite atribuir a diferença exclusivamente a uma delas.

A RSS média variou aproximadamente de 96 a 120 MiB. Essa medida inclui Python e caches do driver no processo compartilhado; não deve ser interpretada como memória exclusiva dos FBOs. Para quatro texturas RGBA8 em 1080p, apenas os anexos de cor intermediários somam cerca de 31,64 MiB, excluindo alvos finais, buffers e overhead do driver.

Os arquivos de amostras, resumos por repetição, capturas e metadados estão preservados na execução \texttt{20261021T030448Z\_7e3fc562}. Gráficos apresentam dispersão e médias; nenhuma diferença foi submetida a teste de significância.